\subsection{Markerless Hand-Based Interaction and Creative Applications}
\label{subsec:sota_interaction}

The subsections so far have evaluated the components that convert a webcam
image into a canonical hand pose.
This final subsection turns from the pipeline to its purpose: the interactive
and creative systems such a pose is built to drive.
The three constraints of \autoref{subsec:gap} apply here as forcefully as they
do to estimation, but they reframe the question.
What matters is no longer how accurately a hand can be reconstructed, but what
input hardware a user must first acquire before the hand can act as a
controller at all.



\subsubsection{Gesture-Based HCI and Mid-Air Interaction}\label{subsubsec:sota_gesture}

The idea of the hand as a direct input channel is one of the oldest in
\acrshort{hci}.
Bolt's \emph{Put-that-there}~\cite{bolt1980putthatthere} coupled spoken commands
with spatial pointing to command shapes on a graphics display, demonstrating
that gesture aided by a second modality gains both naturalness and referential
precision.
The system instrumented the hand with a position sensor; its lasting
contribution was conceptual, establishing the pointing hand as a spatial
reference for an interface.

\begin{figure}[H]
    \centering
    \captionsetup{justification=centering}
    \includegraphics[width=0.8\linewidth, frame]{img/put-that-there-original-sma.jpg}
    \caption{A user pointing at shapes on the Media Room display.
             Adapted from Bolt's \emph{Put-that-there}~\cite{bolt1980putthatthere}.}
    \label{fig:bolt_put_that_there}
\end{figure}

Vision offered a route to the same naturalness without instrumenting the user.
Freeman et al.~\cite{freeman1998vision} argued that a camera can sense pointing
and gesture directly, and that such unencumbered interaction makes computers
easier to use; they demonstrated several vision-controlled graphics
applications, including a game, driven from a single camera.
The maturation of this paradigm is surveyed by Rautaray and
Agrawal~\cite{rautaray2015vision}, who decompose gesture recognition into
detection, tracking, and recognition, and position hand gestures as a natural
alternative to the mouse and keyboard.
Across this line of work the recurring promise is interaction in which the user
holds and wears nothing; the recurring difficulty, until recently, was that a
single RGB camera could not recover a pose precise enough to drive
more than coarse cursor or classification tasks.
The model-free and model-based estimators of \autoref{subsec:sota_estimation}
are what closed that gap, making the hand viable as a fine-grained spatial
controller from commodity input.



\subsubsection{Hand Tracking in Games and Immersive Environments}\label{subsubsec:sota_immersive}

As tracking matured, hand input moved from research demonstrations into shipping
interactive products.
The most capable of these systems, however, reach their accuracy by adding
dedicated sensing hardware.
Weichert et al.~\cite{weichert2013leap} characterise the Leap Motion Controller,
a desktop optical device marketed for hand-gesture-controlled interfaces with a
declared sub-millimetre accuracy; their measurements place the tracking
deviation at roughly 0.2\,mm in static and 1.2\,mm in dynamic conditions.
The device delivers tracking precise enough for gestural game and desktop
control, but it is a separate peripheral the user must purchase and position.

Egocentric tracking integrated into head-mounted displays represents the
current production standard.
Han et al.~\cite{han2020megatrack} present MEgATrack, a real-time articulated
hand-tracking system that runs at 60\,Hz on a PC and 30\,Hz on
a mobile processor and is the default hand-tracking feature on the Oculus Quest
headset, where it powers both input and social presence.
It demonstrates that controller-free hand input is robust enough to serve as a
primary interaction modality in immersive environments, but it recovers pose
from four fisheye monochrome cameras built into the headset rather than from any
commodity sensor.
Both systems establish that hand tracking is mature enough to drive games and
immersive applications directly.
Both also illustrate the cost the present work seeks to avoid: each presupposes
a sensing apparatus---an optical peripheral or a head-mounted multi-camera
rig---that lies beyond the single webcam fixed by \autoref{subsec:gap}.



\subsubsection{Immersive 3D Creation and the Controller Assumption}\label{subsubsec:sota_creative}

The richest application of spatial hand input is three-dimensional content
creation, where the hand becomes a tool for drawing, sculpting, and painting
directly in space.
The paradigm originates with CavePainting~\cite{keefe2001cavepainting}, which
treats a brush stroke as a three-dimensional analogue of its
two-dimensional counterpart and lets an artist create works inside a fully
immersive CAVE using physical props and full-body gesture.
Keefe et al.\ report that artists valued precisely this embodied, whole-body
character of the act of making.
The approach has since matured into widely adopted commercial tools, of which
Google Tilt Brush and Gravity Sketch are the most prominent, in which marks are
deposited in space by moving the hand through it.

The defining feature of these tools, for the purposes of this thesis, is the
input apparatus they assume.
Drey et al.~\cite{drey2020vrsketchin} observe that the available
\acrshort{vr} sketching tools, Tilt Brush and Gravity Sketch among them,
concentrate almost exclusively on unconstrained mid-air drawing, and attribute
this in part to the absence of any input device in \acrshort{vr} other
than the hand-held controller, a device optimised for entertainment rather than
precise creation.
Mid-air drawing in \acrshort{vr} is itself imprecise regardless of the input device:
Arora et al.~\cite{arora2017sketching} show through controlled study
that the absence of a physical drawing surface is a major source of inaccuracy
in \acrshort{vr} sketching.
The state of the art in immersive creation is therefore doubly hardware-bound.
It presupposes a head-mounted display together with a pair of controllers, and
even with that apparatus in hand it trades away the surface-supported precision
of conventional drawing.

This is the gap the creative application of the present system targets.
The painting tool introduced in~\autoref{subsec:contributions} provides
three-dimensional drawing with point-cloud export from the hand alone: no
head-mounted display and no controllers, the hand sensed by the same commodity
webcam and \acrshort{cpu}-only pipeline used throughout, and the canonical
\acrshort{mano} pose recovered upstream
(\autoref{subsec:sota_estimation},~\autoref{subsec:sota_ik}) supplying the
spatial input a controller would otherwise carry.
The dedicated rigs of Tilt Brush and Gravity Sketch achieve a spatial precision
and tracking volume a single RGB camera cannot match; what the present tool
offers in their place is a three-dimensional creative interface that costs the
user no hardware beyond a webcam they already own.
